% Part: methods
% Chapter: proofs
% Section: example-2

\documentclass[../../../include/open-logic-section]{subfiles}

\begin{document}

\olfileid[pt]{mth}{prf}{ex2}

\olsection{Outro Exemplo} 

\begin{prop}
Se $A \subseteq C$, então $A \cup (C \setminus A) = C$.
\end{prop}  

\begin{proof}
  Suponha que $A \subseteq C$.  Queremos mostrar que $A \cup (C
  \setminus A) = C$.
  \begin{quote}
    Começamos observando que esta é uma afirmação condicional. Ela está
    tacitamente quantificada universalmente: a proposição vale para todos os conjuntos
    $A$ e $C$. Assim, $A$ e $C$ são variáveis para conjuntos arbitrários. Para
    provar tal afirmação, supomos o antecedente e provamos o
    consequente.

    Continuamos usando a suposição de que $A \subseteq C$. Vamos
    desdobrar a definição de~$\subseteq$: a suposição significa que
    todos os !!{element}s de~$A$ são também !!{element}s de~$C$. Vamos anotar
    isto---é um fato importante que usaremos ao longo da
    prova.
  \end{quote}
  Pela definição de~$\subseteq$, como $A \subseteq C$, para todo
  $z$, se $z \in A$, então $z \in C$.
  \begin{quote}
    Desdobramos todas as definições que nos são dadas na
    suposição. Agora podemos passar à conclusão. Queremos mostrar
    que $A \cup (C \setminus A) = C$, e por isso montamos uma prova
    de modo semelhante ao último exemplo: mostramos que todo !!{element} de
    $A \cup (C \setminus A)$ é também !!a{element} de~$C$ e,
    reciprocamente, todo !!{element} de $C$ é !!a{element} de $A \cup (C
    \setminus A)$. Podemos abreviar isso para: $A \cup (C \setminus A)
    \subseteq C$ e $C \subseteq A \cup (C \setminus A)$. (Aqui estamos
    fazendo o oposto de desdobrar uma definição, mas isso torna a
    prova um pouco mais fácil de ler.)  Como isso é uma conjunção, temos
    que provar ambas as partes. Para mostrar a primeira parte, isto é, que todo
    !!{element} de $A \cup (C \setminus A)$ é também !!a{element}
    de~$C$, supomos que $z \in A \cup (C \setminus A)$ para um
    $z$~arbitrário e mostramos que $z \in C$. Pela definição de
    $\cup$, podemos concluir que $z \in A$ ou $z \in C \setminus A$
    a partir de $z \in A \cup (C \setminus A)$. Você já deve estar pegando o
    jeito disso.
    \end{quote}
  $A \cup (C \setminus A) = C$ sse $A \cup (C \setminus A) \subseteq
  C$ e $C \subseteq (A \cup (C \setminus A)$.  Primeiro provamos que
  $A \cup (C \setminus A) \subseteq C$.  Seja $z \in A \cup (C
  \setminus A)$. Então, ou $z \in A$ ou $z \in (C \setminus A)$.
  \begin{quote}
    Chegamos a uma disjunção, e a partir dela queremos provar que
    $z \in C$. Fazemos isso usando a prova por casos.
  \end{quote}
  Caso 1: $z \in A$. Como, para todo $z$, se $z \in A$, $z \in C$,
  temos que $z \in C$.
  \begin{quote}
    Aqui usamos o fato registrado anteriormente, que decorreu da
    hipótese da proposição de que $A \subseteq C$.  O primeiro
    caso está completo, e passamos ao segundo caso, $z \in (C
    \setminus A)$.  Lembre-se de que $C \setminus A$ denota a
    \emph{diferença} dos dois conjuntos, isto é, o conjunto de todos os
    !!{element}s de~$C$ que não são !!{element}s de~$A$.  Mas qualquer
    !!{element} de $C$ que não esteja em~$A$ é, em particular, !!a{element} de~$C$.
  \end{quote}
  Caso 2: $z \in (C \setminus A)$.  Isso significa que $z \in C$ e $z
  \notin A$. Então, em particular, $z \in C$.
  \begin{quote}
    Ótimo, provamos a primeira direção. Agora a segunda
    direção. Aqui provamos que $C \subseteq A \cup (C \setminus
    A)$.  Então supomos que $z \in C$ e provamos que $z \in A \cup (C
    \setminus A)$.
  \end{quote}
  Agora seja $z \in C$. Queremos mostrar que $z \in A$ ou $z \in C
  \setminus A$.
  \begin{quote}
    Como todos os !!{element}s de $A$ são também !!{element}s de $C$, e $C
    \setminus A$ é o conjunto de todas as coisas que são !!{element}s de $C$
    mas não de $A$, segue-se que $z$ está ou em $A$ ou em $C
    \setminus A$.  Isso pode ser um pouco obscuro se você ainda não souber
    por que o resultado é verdadeiro.  Seria melhor prová-lo
    passo a passo.  Ajudará usar um fato simples que podemos
    enunciar sem prova: $z \in A$ ou $z \notin A$. Isso é chamado de
    ``princípio do terceiro excluído:'' para qualquer afirmação~$p$, ou
    $p$ é verdadeira ou sua negação é verdadeira. (Aqui, $p$ é a afirmação
    de que $z \in A$.)  Como isso é uma disjunção, podemos novamente usar
    a prova por casos.
  \end{quote}
  Ou $z \in A$ ou $z \notin A$. No primeiro caso, $z \in A \cup
  (C \setminus A)$. No segundo caso, $z \in C$ e $z \notin A$, de modo que
  $z \in C \setminus A$.  Mas então $z \in A \cup (C \setminus A)$.
  \begin{quote}
    Nossa prova está completa: mostramos que $A \cup (C \setminus A) = C$.
  \end{quote}
\end{proof}

\end{document}
